\section{信道与信道容量：一条线最多能传多少信息}

\begin{idea}{这一讲只回答一个问题}
前几讲算的是“一条消息本身值多少 bit”。这一讲把消息真的送出去：\textbf{经过一条会出错的线之后，最多还能有多少 bit 到对面？}
这个“最多能到多少”就叫\textbf{信道容量}，记作 $C$。一句话分工：前几讲管“消息值多少”，这一讲管“\textbf{线只肯给你多少}”。
\end{idea}

你在电话里给对面报一串 0/1。你说“0”，对面大多数时候记下 $0$，但信号一抖，他也可能记成 $1$。
他记成什么你事前不知道，只知道\textbf{各种结果的概率}——这就是信道：不是理想的铜线，而是一句“发了这个，收到那个的可能性多大”。

\subsection{信道就是一张“发什么、收到什么”的概率表}

符号先讲清（每个符号第一次出现我都用大白话说一遍）：
\begin{itemize}[leftmargin=1.6em,itemsep=1pt]
  \item $X$：我\textbf{发出去}的符号，是随机变量（这一格发什么，事前我只知道概率）。
  \item $Y$：对面\textbf{收回来}的符号，也是随机变量。
  \item $\p{y|x}$：读作“\textbf{在已经发了 $x$ 这个前提下，收到 $y$ 的概率}”。那个竖线是中文“在……的条件下”，不是除法、也不是“或者”。
\end{itemize}

\begin{calc}{先学会读这张表（本讲的基本功）}
某条线，输入 $\{0,1\}$，输出也是 $\{0,1\}$，转移概率表是：
\begin{center}
\begin{tabular}{ccc}
\toprule
我发的 $x$ & 收到 $y=0$ 的概率 & 收到 $y=1$ 的概率 \\
\midrule
$x=0$ & $0.9$ & $0.1$ \\
$x=1$ & $0.2$ & $0.8$ \\
\bottomrule
\end{tabular}
\end{center}

\textbf{三条读表规则（务必背下来）：}
\begin{enumerate}[leftmargin=1.8em,itemsep=2pt]
  \item \textbf{横着读一行}：第一行读作“我发 $0$，对面收到 $0$ 的概率 $0.9$、收到 $1$ 的概率 $0.1$”；
        第二行读作“我发 $1$，对面收到 $0$ 的概率 $0.2$、收到 $1$ 的概率 $0.8$”。
  \item \textbf{每行加起来必须等于 $1$}：$0.9+0.1=1$，$0.2+0.8=1$。道理很土——发出去一个符号，对面总得收到点东西，所有可能铺开就是 $100\%$。
  \item \textbf{每列加起来不必等于 $1$}：第一列 $0.9+0.2=1.1$，第二列 $0.1+0.8=0.9$。因为“收到 $0$ 的总概率”取决于我发 $0$、发 $1$ 各占多少，而那是\textbf{我}定的，不是线定的。
\end{enumerate}

写成矩阵：
\[
Q=\begin{pmatrix}
\p{y=0|x=0} & \p{y=1|x=0}\\
\p{y=0|x=1} & \p{y=1|x=1}
\end{pmatrix}
=\begin{pmatrix} 0.9 & 0.1\\ 0.2 & 0.8\end{pmatrix}.
\]
人话：\textbf{一行 = 一种“我发了什么”，一格 = 一个条件概率}；第 $i$ 行第 $j$ 列就是“发了 $i$ 收到 $j$”的概率。
\end{calc}

\begin{defbox}{信道（本课只考“离散无记忆信道”）}
给定输入字母表 $A$、输出字母表 $B$，把所有 $\p{y|x}$ 列成的那张表（记 $Q=(q(y|x))$）就完整描述了一个信道，它叫\textbf{转移概率矩阵}。
“无记忆”是说：这一格出不出错与上一格无关，所以整串符号的概率可以直接连乘，我们只需要研究“一格”的信道。
\end{defbox}

\subsubsection{从这张表走到 $\I{X;Y}$：只有两步}

你已经知道“互信息 = 知道 $Y$ 之后，$X$ 的不确定性减少了多少”。写成公式：
\[
\I{X;Y}=\Hx{Y}-\Hx{Y|X}.
\]
人话：\textbf{看着 $Y$ 我还糊涂多少}，\textbf{减去}“本来就知道会糊涂多少”，剩下的才是真传过来的。
其中 $\Hx{Y|X}$ 是“在知道发了什么的条件下，对收到什么的平均不确定性”，算法是\textbf{每行各自算熵，再按我发各符号的概率加权平均}：
\[
\Hx{Y|X}=\sum_{x}\p{x}\,\Hx{Y|X=x}
=\sum_{x}\p{x}\left[-\sum_{y}\p{y|x}\lb\p{y|x}\right].
\]
人话：$\Hx{Y|X}$ 就是“\textbf{这条线平均每格给你添多少乱}”。它是线自己的毛病，与你怎么发无关——这一点后面要用。

\begin{calc}{手算示范：上面那张表，输入等概，求 $\I{X;Y}$}
设 $\p{X=0}=\p{X=1}=0.5$。（约定 $\Hx{p}=-p\lb p-(1-p)\lb(1-p)$，读作“一个概率为 $p$ 的二元事件的熵”，单位 bit；故 $\Hx{0.5}=1$，$\Hx{0}=\Hx{1}=0$。）

\textbf{第 1 步，全概率求输出的分布}：$\p{Y=0}=0.5\times0.9+0.5\times0.2=0.45+0.10=0.55$，$\p{Y=1}=1-0.55=0.45$。

\textbf{第 2 步，两处熵}。用换底 $\lb n=\dfrac{\ln n}{\ln2}$，$\ln2=0.6931$，查得 $\lb0.55=-0.8625$、$\lb0.45=-1.1520$、$\lb0.9=-0.1520$、$\lb0.2=-2.3219$、$\lb0.8=-0.3219$，于是
\[
\Hx{Y}=0.55\times0.8625+0.45\times1.1520=0.4744+0.5184=0.9928\bit,
\]
\[
\Hx{Y|X}=0.5\times\Hx{0.1}+0.5\times\Hx{0.2}=0.5\times0.4690+0.5\times0.7219=0.5954\bit .
\]
（其中 $\Hx{0.1}=0.1\times3.3219+0.9\times0.1520=0.4690$，$\Hx{0.2}=0.2\times2.3219+0.8\times0.3219=0.7219$。）

\textbf{第 3 步，相减}：$\I{X;Y}=0.9928-0.5954=0.3974\bit$。
人话：每传一格平均只有 $0.397$ bit 真到了对面，而理论上限是 $\Hx{X}=\lb2=1$ bit，线吃掉了大头。
\end{calc}

\begin{idea}{记住这个“手算三步走”}
\textbf{第一步}全概率算 $\p{Y=y}$；\textbf{第二步}每行求熵再加权得 $\Hx{Y|X}$，顺带算 $\Hx{Y}$；\textbf{第三步}相减。最容易错的是第一步忘了乘 $\p{x}$。
\end{idea}

\subsection{三个具体信道：把 $\I{X;Y}$ 一个一个算出来}

\subsubsection{无噪信道：收到就等于知道发了什么}

发 $0$ 必收 $0$，发 $1$ 必收 $1$，表是

\begin{center}
\begin{tabular}{ccc}
\toprule
我发的 $x$ & 收到 $y=0$ & 收到 $y=1$ \\
\midrule
$x=0$ & $1$ & $0$ \\
$x=1$ & $0$ & $1$ \\
\bottomrule
\end{tabular}
\end{center}

每行都是“一个 $1$ 一个 $0$”，每行熵 $=\Hx{1}=-1\cdot\lb1-0\cdot\lb0=0$（约定 $0\lb0=0$），故 $\Hx{Y|X}=0$，
\[
\I{X;Y}=\Hx{Y}-0=\Hx{Y}=\Hx{X},
\qquad
C=\max_{\p{x}}\Hx{X}=\lb2=1\bit .
\]
人话：\textbf{线一点乱都不添，发多少 bit 就收到多少 bit}；等概时 $X$ 最“不确定”（$\Hx{X}=1$ bit），所以容量是 $1$ bit。
换成 $4$ 个符号的无噪信道（$00,01,10,11$ 都不错），等概时 $\Hx{X}=\lb4=2\bit$，故 $C=2$ bit。
规律：\textbf{无噪信道 $C=\lb|X|$}（$|X|$ 是输入符号个数）。

\subsubsection{二元对称信道 BSC：以概率 $p$ 翻错}

“二元”= 发 $0/1$ 收 $0/1$；“对称”= 两种错误机会一样，都以概率 $p$ 翻错。

\begin{center}
\begin{tabular}{ccc}
\toprule
我发的 $x$ & 收到 $y=0$ & 收到 $y=1$ \\
\midrule
$x=0$ & $1-p$ & $p$ \\
$x=1$ & $p$ & $1-p$ \\
\bottomrule
\end{tabular}
\end{center}

两行是\textbf{同一组数 $(1-p,\,p)$ 换了顺序}——这个性质到“三个必背结论”那一节会变成一条能直接套的公式。
设 $\p{X=0}=p_0$，则 $\p{Y=0}=p_0(1-p)+(1-p_0)p=p+p_0(1-2p)$（这行只是告诉你输出分布跟 $p_0$ 有关）。
而两行熵相同，加权平均还是它自己：
\[
\Hx{Y|X}=p_0\Hx{p}+(1-p_0)\Hx{p}=\Hx{p}
=-\big[p\lb p+(1-p)\lb(1-p)\big].
\]
人话：\textbf{这条线每格固定吃掉 $\Hx{p}$ bit}，不管你发得均不均匀。于是
\[
\I{X;Y}=\Hx{Y}-\Hx{p}\le\lb2-\Hx{p}=1-\Hx{p},
\]
因为 $Y$ 只有两个取值，熵最大是 $\lb2=1$。什么时候取等？当 $\Hx{Y}=1$，即 $Y$ 等概；取 $p_0=0.5$ 时 $\p{Y=0}=p+0.5(1-2p)=0.5$，等号能取到。所以
\[
\boxed{\ C=\max_{p_0}\big[\Hx{Y}-\Hx{p}\big]=1-\Hx{p}\ }
\]
人话：容量 = 天花板 $1$ bit 减去“翻错带来的糊涂”$\Hx{p}$，最佳输入是等概。

\begin{calc}{代入 $p=0.1$：$C\approx0.531$ bit（算术一步不跳）}
\[
\Hx{0.1}=-\big[\,0.1\lb0.1+0.9\lb0.9\,\big],
\]
先查对数：$\lb0.1=\dfrac{\ln0.1}{\ln2}=\dfrac{-2.3026}{0.6931}=-3.3219$，$\lb0.9=\dfrac{\ln0.9}{\ln2}=\dfrac{-0.1054}{0.6931}=-0.1520$。代入：
\[
\Hx{0.1}=0.1\times3.3219+0.9\times0.1520=0.3322+0.1368=0.4690\bit,
\]
\[
C=1-\Hx{0.1}=1-0.4690=0.5310\bit\approx0.531\bit .
\]
人话：一条“错 10\%”的线，每格平均只送到 $0.531$ bit，近一半被线吃掉。这个 $0.469$ 请记住，考试很可能要你写。

\textbf{两个极端}：$p=0$ 时 $\Hx{0}=-0\lb0-1\lb1=0$，$C=1-0=1\bit$（退化成无噪信道）；
$p=0.5$ 时 $\Hx{0.5}=1$，$C=1-1=0\bit$。中间再记两个点：$p=0.01$ 时 $\Hx{0.01}=0.0808$，$C=0.919\bit$；
$p=0.11$ 时 $\Hx{0.11}=0.4999\approx0.5$，$C\approx0.500\bit$——\textbf{错误率 11\% 就把容量砍掉一半}。
\end{calc}

\subsubsection{二元擦除信道 BEC：以概率 $p$ 变成“擦除符号 ?”}

这条线不偷偷改，而是\textbf{大声告诉你“这位我丢了”}。输出有三个符号：$0$、$1$、擦除记号 $?$。

\begin{center}
\begin{tabular}{cccc}
\toprule
我发的 $x$ & 收到 $y=0$ & 收到 $y=?$ & 收到 $y=1$ \\
\midrule
$x=0$ & $1-p$ & $p$ & $0$ \\
$x=1$ & $0$ & $p$ & $1-p$ \\
\bottomrule
\end{tabular}
\end{center}

每行是 $(1-p,p,0)$ 或 $(0,p,1-p)$，熵都是 $\Hx{p}$，所以 $\Hx{Y|X}=\Hx{p}$。取 $\p{X=0}=0.5$，则
\[
\p{Y=?}=0.5p+0.5p=p,\qquad \p{Y=0}=\p{Y=1}=\frac{1-p}{2},
\]
\[
\Hx{Y}=\underbrace{\Hx{p}}_{\text{是不是被擦}}+\underbrace{(1-p)\cdot\lb2}_{\text{没被擦时是 }0\text{ 还是 }1}
=\Hx{p}+(1-p).
\]
（第二步用了熵的链式展开：先问“擦没擦”，没擦的那部分占 $1-p$、里面 $0/1$ 各半、熵为 $\lb2=1$ bit，再乘 $1-p$。）相减：
\[
\I{X;Y}=\big[\Hx{p}+(1-p)\big]-\Hx{p}=(1-p)\bit,
\qquad
\boxed{\ C=1-p\ }
\]
人话：\textbf{擦掉 $p$ 的比例，就正好损失 $p$ 的容量}，一点不冤。$p=0.1$ 时 $C=0.9\bit$；$p=0.5$ 时 $C=0.5\bit$；$p=1$ 时 $C=0$。

\begin{idea}{BSC 和 BEC 都“错 50\%”，容量为什么差这么多？}
BSC 的 $C=0$，BEC 的 $C=0.5\bit$。差别在于\textbf{擦除是诚实的错误，翻错是撒谎}：擦掉的那位你明知没用可以扔掉，
剩下 $1-p$ 的比例仍敢用；而翻错 50\% 时你连“哪一位错了”都不知道，收到 $0$ 时发 $0$ 和发 $1$ 一样可能，一格信息都没有。
\end{idea}

\subsection{信道容量的定义：我能选的只有“发什么符号的概率”}

上面三个例子里 $\I{X;Y}$ 都随“我怎么发”变化。为了把线本身的好坏单独拎出来：
\begin{defbox}{信道容量}
\[
C=\max_{\p{x}}\I{X;Y}
\]
人话：\textbf{在“发各符号的概率分布”这个我能自由选择的旋钮上，挑一个让传过去的信息最大的，那个最大值就是这条线的容量。}
\end{defbox}

三个必须分清的点：
\begin{enumerate}[leftmargin=1.8em,itemsep=2pt]
  \item \textbf{我能改的只有 $\p{x}$。}转移概率 $\p{y|x}$ 是线本身的性质，我改不了（要改就得换线），所以“取最大”只能对 $\p{x}$ 取。
  \item \textbf{$C$ 是线的属性，与我的输入分布无关。}同一条线，我发得聪明就能传到 $C$、发得笨就传得少，但这条线的 $C$ 一直是那个数。
        像一条公路的“最高时速”是路和车的能力，不取决于你今天开多快。
  \item \textbf{求 $C$ 时又要对 $\p{x}$ 取最大。}两句不矛盾：$C$ \textbf{是}一个不依赖 $\p{x}$ 的数，但它的\textbf{算法}是对 $\p{x}$ 求最大值。考试最爱在这儿挖坑。
\end{enumerate}

\begin{defbox}{最佳输入分布（考试会直接问这个词）}
使 $\I{X;Y}$ 取到最大值 $C$ 的那个输入分布 $\p{x}$，叫\textbf{最佳输入分布}。
$\I{X;Y}$ 关于 $\p{x}$ 是上凸的（中间高两头低，不会有两个一样高的山头），所以最佳分布只有一个。
\end{defbox}

\begin{defbox}{判别最佳分布的条件（好用，且能当验算）}
记 $I(X=k;Y)$ 为“\textbf{固定发第 $k$ 个符号时}这一格能传过去的互信息”（把“$x=k$ 板上钉钉”代入 $\I{X;Y}$ 就能算）。则最佳分布 $\p{x}^*$ 满足：
凡 $\p{x=k}>0$（用得上）的，$I(X=k;Y)=C$；凡 $\p{x=k}=0$（不用）的，$I(X=k;Y)\le C$。
\end{defbox}

人话：\textbf{最佳方案让“每个被用到的符号都同样划算”}，不划算的符号干脆不给概率。
它有两个用处：一是反手验算（算出 $p^*$ 后代回去看是否都等于 $C$），二是找 $p^*$（让各 $I(X=k;Y)$ 相等，解方程）——后面“非对称信道”那一节就这么干。

\subsection{三个必背结论}

\subsubsection{结论一：$C\le\lb|X|$（容量有天花板）}

推理一行：$\I{X;Y}\le\Hx{X}$（知道 $Y$ 能让 $X$ 的不确定性减少，但不会减少成负的、也不会比原来更多），而 $\Hx{X}\le\lb|X|$（$|X|$ 个符号的熵最大是等概的 $\lb|X|$）。所以
\[
\boxed{\ C\le\lb|X|\ }
\qquad
\text{人话：一格最多传 $\lb|X|$ bit，因为这一格我最多只能“有 $\lb|X|$ bit 的疑问”。}
\]
\textbf{什么时候取等}：$\I{X;Y}=\Hx{X}$，即 $\Hx{X|Y}=0$（收到 $Y$ 就完全确定 $X$），也就是\textbf{无噪信道}且输入等概。
\textbf{数值}：输入 $3$ 个符号则 $C\le\lb3=1.585\bit$；$4$ 个符号则 $C\le\lb4=2\bit$。
\textbf{常常够不着的例子}（课本例 5.1.2“复制信道”）：$26$ 个字母，每个以 $0.5$ 保持原样、以 $0.5$ 变成下一个字母，则 $C=\lb13=3.700\bit$，
而天花板是 $\lb26=4.700\bit$，少了整整 $1$ bit——正好是“到底有没有错一步”这个二元问题的 $1$ bit。

\subsubsection{结论二：强对称（对称）信道 $C=\lb|X|-\Hx{\text{一行}}$}

\textbf{什么叫对称}：转移概率矩阵的\textbf{每一行都是第一行换了顺序}（行可排列），并且\textbf{每一列元素之和相同}。二元对称信道、$K$ 元对称信道都属于这类。

为什么公式长这样：行可排列 $\Rightarrow$ 各行熵相同 $\Rightarrow$ $\Hx{Y|X}=\Hx{\text{一行}}$，与输入分布无关；
列和相等 $\Rightarrow$ 输入等概时输出也等概 $\Rightarrow$ $\Hx{Y}$ 顶到天花板 $\lb|X|$（这类信道输入输出符号数相同）。于是
\[
\boxed{\ C=\lb|X|-\Hx{\text{一行}},\qquad \text{最佳输入分布 = 等概分布}\ }
\]
人话：\textbf{容量 = 天花板减掉“一行那点平均糊涂”}，$\Hx{\text{一行}}$ 就是这条线平均每格自己添的乱。

\textbf{用它重算 BSC，检验公式靠不靠谱}：BSC 每行都是 $(1-p,p)$，这一行的熵就是 $\Hx{p}$，故 $C=\lb2-\Hx{p}=1-\Hx{p}$，
与上面 BSC 硬算的结果\textbf{完全一样}；$p=0.1$ 时 $C=1-0.4690=0.531\bit$ 也一样。公式和手算对上过一次，以后就能放心套。

\textbf{再看一个三符号的例子。}某 $3$ 元对称信道：发对的概率是 $0.7$，错成另外两个符号各 $0.15$，
因此它的每一行都是 $(0.7,\;0.15,\;0.15)$。由 $\lb0.7=-0.5146$、$\lb0.15=-2.7370$ 得
\[
\Hx{\text{一行}}=0.7\times0.5146+2\times0.15\times2.7370=0.3602+0.8211=1.1813\bit,
\]
\[
C=\lb3-1.1813=1.5850-1.1813=0.4037\bit ,
\]
最佳输入是三个符号各 $1/3$。注意它远小于 $1$ bit——三个符号看着热闹，其实穿不过 $0.41$ bit。

\subsubsection{结论三：准对称信道，等概输入也最佳}

\textbf{什么叫准对称}：行可排列（每行都是第一行换顺序），但\textbf{列和不相等}。此时输出符号可分成几组 $B_1,\dots,B_S$，使每组对应的列组成的子矩阵仍对称。容量公式：
\[
C=\sum_{s=1}^{S}N_s\lb\frac{|B_s|}{N_s}-\Hx{\text{一行}},
\qquad N_s=\text{一行里落进第 }s\text{ 组的概率（对任何输入都相同）}.
\]
人话：\textbf{把输出分成几个“小组”，每组的容量照对称公式算，再按概率加起来}。最佳输入仍是等概（题目只要求你会用这一条）。

\textbf{例子}（课本例 5.1.7，既有擦除又有误传）：每行是 $(1-p-q,\ q,\ p)$，取 $p=q=0.1$。
分两组：$B_1=\{0,1\}$ 是“收到真符号”（$|B_1|=2$，$N_1=(1-p-q)+p=1-q=0.9$），$B_2=\{?\}$ 是“擦除”（$|B_2|=1$，$N_2=q=0.1$），
$N_1+N_2=1$。代入，其中 $\lb\dfrac{2}{0.9}=\lb2.2222=1.1520$、$\lb\dfrac{1}{0.1}=\lb10=3.3219$、
$\Hx{0.8,0.1,0.1}=0.8\times0.3219+2\times0.1\times3.3219=0.2575+0.6644=0.9219$：
\[
C=0.9\times1.1520+0.1\times3.3219-0.9219=1.0368+0.3322-0.9219=0.4471\bit .
\]
\textbf{验算}（输入等概，直接算 $\Hx{Y}-\Hx{Y|X}$）：$\p{Y=0}=\p{Y=1}=\dfrac{1-p}{2}=0.45$、$\p{Y=?}=p=0.1$，故
$\Hx{Y}=1.3690$、$\Hx{Y|X}=0.9219$、$\I{X;Y}=0.4471\bit$，与公式结果一样，说明公式没记错。
\textbf{考试拿不准就用这招验算。}

\subsection{非对称信道：套不了公式的时候怎么办}

把最常见的坑摆出来。考虑 $2\times2$ 信道
\[
Q=\begin{pmatrix}1-\alpha & \alpha\\ \beta & 1-\beta\end{pmatrix}
\quad\text{取 }\alpha=0.1,\ \beta=0.3,\qquad
Q=\begin{pmatrix}0.9 & 0.1\\ 0.3 & 0.7\end{pmatrix},
\]
第一行是“发 $0$：$90\%$ 收对、$10\%$ 收成 $1$”，第二行是“发 $1$：$30\%$ 收成 $0$、$70\%$ 收对”。

\textbf{第一件事：判对称性。}两行 $(0.9,0.1)$ 与 $(0.3,0.7)$ \textbf{不是同一组数的重新排列}（数都不一样），列和 $0.9+0.3=1.2\ne0.1+0.7=0.8$，
所以它\textbf{既不对称也不准对称，前面的公式一个都不能用}。这就是“先判对称性”的意思。

\textbf{第二件事：把 $\I{X;Y}$ 写成一个只含 $p$ 的函数}，其中 $p=\p{X=0}$，$1-p=\p{X=1}$。每行熵是常数：
由 $\lb0.3=-1.7370$、$\lb0.7=-0.5146$ 得 $\Hx{0.3}=0.3\times1.7370+0.7\times0.5146=0.5211+0.3602=0.8813$，故
\[
\Hx{Y|X}=0.4690p+0.8813(1-p).
\]
输出分布\textbf{这次随 $p$ 变}（小心）：
\[
q:=\p{Y=0}=0.9p+0.3(1-p)=0.3+0.6p,
\qquad
\p{Y=1}=1-q=0.7-0.6p .
\]
于是
\[
\I{X;Y}=\Hx{q}-0.4690p-0.8813(1-p)
=\Hx{0.3+0.6p}-0.4690p-0.8813(1-p).
\]
人话：一个自变量、一条曲线，最大值只可能在两端或导数为 $0$ 处。

\textbf{第三件事：先扫描看大概。}（考试时间紧时，扫描表本身就是分。）

\begin{center}
\begin{tabular}{cccccc}
\toprule
$p$ & $0$ & $0.2$ & $0.4$ & $0.5$ & $0.8$ \\
\midrule
$q=0.3+0.6p$ & $0.30$ & $0.42$ & $0.54$ & $0.60$ & $0.78$ \\
$\Hx{q}$ & $0.8813$ & $0.9815$ & $0.9954$ & $0.9710$ & $0.7602$ \\
$\Hx{Y|X}$ & $0.8813$ & $0.7988$ & $0.7164$ & $0.6751$ & $0.5515$ \\
$\I{X;Y}$ & $0$ & $0.1826$ & $0.2790$ & $0.2958$ & $0.2087$ \\
\bottomrule
\end{tabular}
\end{center}

（$\Hx{Y|X}$ 那一行举例：$p=0.2$ 时 $0.4690\times0.2+0.8813\times0.8=0.0938+0.7050=0.7988$。表中数字四舍五入到四位，故末行与前几行相减可能差 $0.0001$。）
看得出：$p=0$ 与 $p=1$ 时 $\I{X;Y}=0$（只发一种符号，等于什么也没说），最大值在 $0.5$ 偏右一点。

\textbf{第四件事：求导定精确值。}只需一个工具：
\[
\frac{\mathrm{d}}{\mathrm{d}u}\Hx{u}=\lb\frac{1-u}{u},
\qquad\text{人话：熵函数在 } u \text{ 处的斜率是 } \lb\frac{1-u}{u}.
\]
（对 $\Hx{u}=-u\lb u-(1-u)\lb(1-u)$ 直接求导，两项合起来正好是上式。）链式法则给出 $\dfrac{\mathrm{d}}{\mathrm{d}p}\Hx{q}=\dfrac{\mathrm{d}q}{\mathrm{d}p}\lb\dfrac{1-q}{q}=0.6\lb\dfrac{1-q}{q}$，于是
\[
\frac{\mathrm{d}\I{X;Y}}{\mathrm{d}p}=0.6\ \lb\frac{1-q}{q}-0.4690+0.8813=0.6\ \lb\frac{1-q}{q}+0.4123 .
\]
令它等于 $0$：$\lb\dfrac{1-q}{q}=-\dfrac{0.4123}{0.6}=-0.6872$，即 $\dfrac{1-q}{q}=2^{-0.6872}=0.6211$。
解出 $1-q=0.6211q\Rightarrow q=\dfrac{1}{1.6211}=0.6169$，回代 $q=0.3+0.6p$：
\[
p^*=\frac{0.6169-0.3}{0.6}=\frac{0.3169}{0.6}=0.5281 .
\]
（二阶导数为负，所以这是最大值点而不是最小值点。）

\textbf{第五件事：算容量}，把 $p^*=0.5281$、$q^*=0.6169$ 代回。先算 $\Hx{0.6169}$：$\lb0.6169=-0.6969$、$\lb0.3831=-1.3842$，故
\[
\Hx{0.6169}=0.6169\times0.6969+0.3831\times1.3842=0.4299+0.5303=0.9602\bit,
\]
再把线吃掉的部分加权：$0.46900\times0.52812=0.24769$（发 $0$ 那行）、$0.88129\times0.47188=0.41586$（发 $1$ 那行），两项和 $0.66355$。相减：
\[
\boxed{\ C=0.96020-0.66355=0.29667\approx0.2967\bit,\qquad p^*\approx0.5281\ }
\]
对照：若硬算等概输入（$p=0.5$），只能得 $\I{X;Y}=0.2958\bit$，比最佳值少约 $0.001\bit$——
本例最优点离等概很近，但\textbf{严格说不是等概}。

\textbf{第六件事：用判别条件验算}（强烈建议写进卷面）。把 $p^*$ 代回，分别算两个符号的 $I(X=k;Y)$：
\begin{align*}
I(X=0;Y)&=0.9\lb\frac{0.9}{0.6169}+0.1\lb\frac{0.1}{0.3831}\\
&=0.9\times0.5449+0.1\times(-1.9378)=0.4904-0.1938=0.2966,
\end{align*}
\begin{align*}
I(X=1;Y)&=0.3\lb\frac{0.3}{0.6169}+0.7\lb\frac{0.7}{0.3831}\\
&=0.3\times(-1.0400)+0.7\times0.8695=-0.3120+0.6086=0.2966,
\end{align*}
两者都等于 $C\approx0.2967$（四位小数下都写成 $0.2966$，差别纯属四舍五入）；两个符号都在使用（$p^*$ 与 $1-p^*$ 都大于 $0$），正好符合判别条件。

\begin{idea}{这个例子要记住的三句话}
（1）\textbf{不对称就不能套公式}，先看两行是不是同一组数的排列；
（2）非对称时\textbf{把 $\I{X;Y}$ 写成只含 $p$ 的一元函数，再对 $p$ 求最大}（求导或扫描）；
（3）答案通常是“$p^*$ 略偏等概、$C$ 明显偏小”——不对称本身就是一种额外损失。
\end{idea}

\subsection{级联信道：串联两道关卡只会更差}

“级联”就是串起来：$X\to$ 信道 1 $\to Y\to$ 信道 2 $\to Z$。真实通信几乎都这样（发端、线路、收端各有瑕疵）。
\textbf{关键事实}：给定 $Y$ 之后，$Z$ 不再额外知道 $X$ 的任何事——$X$ 的信息只能“经 $Y$ 转手”。于是
\[
\I{X;Z}\le\I{X;Y},
\qquad
\text{人话：从终点 } Z \text{ 看到的 } X,\text{ 不会比中间站 } Y \text{ 看到的更多。}
\]
这就是有名的\textbf{数据处理不等式}：任何“后处理”（复印、排序、画图、压缩）都不能凭空造出信息。两条不等式一夹：
\[
\boxed{\ C_{\text{级联}}\le\min\{C_1,\ C_2\}\ }
\qquad
\text{人话：串起来不超过任何一环，瓶颈在最短的那块板。}
\]
证明只需两行：$\I{X;Z}\le\I{X;Y}\le C_1$，同时 $\I{X;Z}\le\I{Y;Z}\le C_2$（$X\to Y\to Z$ 也构成马尔可夫链），
故 $\I{X;Z}\le\min\{C_1,C_2\}$；对输入分布取最大，右边是常数，不等号保住。

\begin{calc}{两个 BSC 级联，手算一次}
单条线都是 BSC，翻错概率 $\varepsilon=0.1$，$C_1=1-\Hx{0.1}=0.531\bit$。串两条后，只有“先对后错”或“先错后对”才净翻错：
\[
\varepsilon_2=0.9\times0.1+0.1\times0.9=0.18,
\qquad\text{一般式 } \varepsilon_2=\frac{1-(1-2\varepsilon)^2}{2}=\frac{1-0.64}{2}=0.18 .
\]
（$(1-2\varepsilon)^2=0.8^2=0.64$ 叫“可靠性”，串级时相乘，所以错误率被拉向 $0.5$。）于是
\[
C_{\text{级联}}=1-\Hx{0.18}=1-0.6801=0.3199\bit<0.531\bit,
\]
其中 $\Hx{0.18}=0.18\lb\dfrac{1}{0.18}+0.82\lb\dfrac{1}{0.82}=0.18\times2.4739+0.82\times0.2863=0.4453+0.2348=0.6801$。
确实更差，符合 $C_{\text{级联}}\le\min\{0.531,0.531\}=0.531$。串 $N$ 条时 $\varepsilon_N=\frac{1-(1-2\varepsilon)^N}{2}\to0.5$，
故 $C_N\to1-\Hx{0.5}=0$——\textbf{无限串联等于完全失联}。
\end{calc}

\begin{idea}{为什么这个结论值得单独记}
它解释了一件反直觉的事：\textbf{多一道环节永远不是好事}。所以工程上能省一环就省一环，
也说明“反复转发、先解密再加密”救不了已经丢掉的信息。
\end{idea}

\subsection{高斯信道：连续世界的容量公式（只记结论和直觉）}

如果信号是连续电压、噪声是热噪声（高斯分布），就是\textbf{高斯信道}。结论漂亮得不像话：
\[
\boxed{\ C=\frac12\lb\!\left(1+\frac{P}{N}\right)\ \text{bit/符号}\ }
\]
$P$ 是信号平均功率（我能花多大力气说话），$N$ 是噪声平均功率（环境多吵），$\dfrac{P}{N}$ 叫\textbf{信噪比}，
读作“\textbf{信号功率比噪声功率}”：$P/N=1000$ 就是信号比噪声强 $1000$ 倍，$P/N=1$ 就是声音和噪音一样大。

\begin{calc}{代入三个数}
\textbf{（i）$P/N=1000$}：$C=\frac12\lb1001$，而 $\lb1001=\dfrac{\ln1001}{\ln2}=\dfrac{6.9088}{0.6931}=9.9672$，故 $C=\frac12\times9.9672=4.9836\approx4.98$ bit/符号。\\
\textbf{（ii）$P/N=1$}：$C=\frac12\lb2=\frac12\times1=0.5$ bit/符号。\\
\textbf{（iii）两个整数好算的}：$P/N=3$ 时 $C=\frac12\lb4=\frac12\times2=1$ bit；$P/N=15$ 时 $C=\frac12\lb16=2$ bit。
\end{calc}

\begin{idea}{为什么是 $\frac12$ 和 $\lb$（一句话直觉）}
\textbf{“为什么是 $\lb$”}：容量只能随信噪比\textbf{对数}增长。看上面的数：$P/N$ 从 $1$ 涨到 $3$（涨 3 倍）多 $0.5$ bit；
从 $3$ 涨到 $15$（涨 5 倍）只多 $1$ bit。\textbf{功率翻一番，只多 0.5 bit}——这就是“\textbf{对数瓶颈}”：靠加功率换速率，代价是指数级的。
\textbf{“为什么有 $\frac12$”}：连续信号取值无穷多，直接算熵没法比；$\frac12$ 是“每自由度”的折算（功率受限时高斯分布的熵最大，代进去系数正好是 $\frac12$）。
本讲只要求记公式、会代入、会解释 $P/N$ 是什么。
\end{idea}

\subsection{应试与答疑}

\begin{exam}{考试怎么考：见到信道题的固定动作}
\textbf{标准框架（按顺序写，一行都别省）：}
\begin{enumerate}[leftmargin=1.8em,itemsep=2pt]
  \item \textbf{抄转移矩阵}，检查每行和是否为 $1$（不是就先怀疑抄错题）。
  \item \textbf{判对称性}（这步决定后面用公式还是列方程）：每行都是第一行换顺序 $\Rightarrow$ 输入对称；
        行可排列\textbf{且}列和相等 $\Rightarrow$ 强对称（可直接套公式）；行可排列但列和不等 $\Rightarrow$ 准对称（等概最佳，按分组公式或直接算 $\Hx{Y}-\Hx{Y|X}$）；
        连行可排列都不是（如 $\left(\begin{smallmatrix}0.9&0.1\\0.3&0.7\end{smallmatrix}\right)$）$\Rightarrow$ \textbf{不许套公式}。
  \item \textbf{对称/准对称}：写 $C=\lb|X|-\Hx{\text{一行}}$（或分组公式）$\to$ 代数 $\to$ 补一句“最佳输入分布为等概分布，即 $p=1/|X|$”。
  \item \textbf{非对称}：设 $p=\p{X=0}$ $\to$ 全概率写 $q=\p{Y=0}$（检查 $\p{Y=1}=1-q$）$\to$ 写 $\I{X;Y}=\Hx{q}-\sum_x\p{x}\Hx{\text{该行}}$
        $\to$ 求导令 $0$（或给扫描表）$\to$ 解 $p^*$ $\to$ 代回得 $C$ $\to$ 用“各 $I(X=k;Y)=C$”验算。
  \item \textbf{级联题}：先写 $C_{\text{级联}}\le\min\{C_1,C_2\}$，再按需把矩阵相乘求等效信道；问为什么就一句“级联是马尔可夫链 + 数据处理不等式”。
  \item \textbf{高斯题}：$C=\frac12\lb(1+P/N)$，单位 bit/符号；先把 $P/N$ 算出来。
\end{enumerate}
\textbf{三个最常见的扣分点}：（1）不判对称性就套公式（本例会虚高到 $0.53$，正确是 $0.297$）；
（2）算 $\Hx{Y|X}$ 忘了乘 $\p{x}$，直接把两行熵相加；（3）说“$C$ 跟输入分布有关”——$C$ 与输入分布无关，只是求它时要对输入分布取最大。
\end{exam}

\begin{qa}{为什么 $C$ 是对“输入分布”取最大，而不是对“转移概率”取最大？}
因为\textbf{转移概率是线的物理性质，我改不了}——它由设备、线路、噪声决定，给定之后这条线的身份就定了；
而 $\p{x}$（我打算按什么比例发 $0$ 和 $1$）是我完全能控制的旋钮。$C$ 要回答“\textbf{这条线最好能给我多少}”，
所以要在“我能动的那一维”上取最好。反过来若对 $\p{y|x}$ 取最大，等于假设我换一条完全不同的线，跟这条线的容量无关；
而且 $\I{X;Y}$ 对 $\p{y|x}$ 是凸函数（不是凹的），最大值只会出现在“极端”信道上（比如无噪），毫无意义。
\end{qa}

\begin{qa}{$p=0.5$ 的 BSC 为什么容量是 $0$？我明明收到了东西啊。}
要分清“收到了符号”和“收到了\textbf{信息}”。$p=0.5$ 时 $\p{y=0|x=0}=\p{y=0|x=1}=0.5$、$\p{y=1|x=0}=\p{y=1|x=1}=0.5$，
即\textbf{知道收到什么，对“发的是什么”毫无提示}（看到 $Y$ 后 $X$ 的不确定性不减少）。公式也给 $C=1-\Hx{0.5}=1-1=0$。
物理上这叫“输出与输入统计独立”：线在替你掷硬币，结果跟你说了什么无关。对比同是“错一半”的 BEC 容量是 $0.5\bit$，
因为它诚实地说“这位丢了”，没丢的那一半你仍敢用。\textbf{“错”不可怕，“错得让你不知道错了”才可怕。}
\end{qa}

\begin{qa}{等概输入是最佳分布，那输出也一定等概吗？}
不一定。\textbf{对称信道}里两者同时成立，所以 $C=\lb|X|-\Hx{\text{一行}}$ 里的“天花板”项能拿到；
\textbf{非对称信道}就未必：本讲的例子最佳输入 $p^*=0.5281$（不是 $0.5$），对应输出 $q^*=0.3+0.6\times0.5281=0.6169$（也不是 $0.5$）。
一句口诀：\textbf{对称看等概，非对称要列方程}。
\end{qa}
